\chapter{Introduction}

% TODO: Expland more on citations + clean up grammar and spellings

\section{Context and Motivation}
In recent years, the economic and operational value of data has surpassed that of
the software systems which process it. Organisations are increasingly treating
data as a product, complete with ownership structures, stewardship
responsibilities, and lifecycle management requirements. This evolution has
coincided with escalating regulatory scrutiny across jurisdictions. Frameworks
such as APRA’s Prudential Standards \cite{cps234,cps230,cps510,australian_prudential_regulation_authority_prudential_2013}
among international instruments such as the EU’s General Data
Protection Regulation (GDPR) and the United States’ HIPAA mandate not only secure the
handling of information assets, but also demand transparent processes that
demonstrate accountability and control.

Extract--Transform--Load (ETL) workflows lie at the centre of this
transformation. They serve as the operational backbone of analytical systems,
ensuring that data is accurately extracted, transformed, and loaded into
governed repositories. As Mahmud and Ikbal~\cite{mahmud_role_2022} emphasise,
modern ETL pipelines have evolved into socio-technical infrastructures that
sustain the scalability, reliability, and legitimacy of data-driven
decision-making. Yet, despite the rise of DataOps practices, many data pipelines
remain opaque, lacking end-to-end lineage, reproducibility, and verifiable
controls. In regulated industries, this opacity undermines the ability to
demonstrate compliance and erodes institutional trust.

The field of software engineering, by contrast, has already confronted similar
challenges. DevSecOps—the integration of development, security, and
operations—has matured into a discipline that embeds verification, automation,
and governance throughout the software delivery lifecycle
\cite{anisetti_devsecops-based_2022,castellanos_accordant_2021,anisetti_semi-automatic_2020}.
If data is now the product, then DevSecOps provides an architectural and
cultural foundation for treating data pipelines with the same rigour
historically applied to software systems. This research explores that proposition
by examining how DevSecOps principles can be adapted to enable auditability and
\emph{data governance} within ETL workflows.

\section{Problem Statement}
Although the DataOps movement has improved automation and observability within
ETL processes, its focus has largely been operational efficiency rather than
regulatory assurance~\cite{munappy_ad-hoc_2020,mahmud_role_2022}. Current
orchestration frameworks, while powerful, rarely produce immutable audit trails,
reproducible execution states, or cryptographically verifiable lineage, leaving
auditability an emergent rather than designed property
\cite{wang_enabling_2011,weigand_auditability_2019}. Consequently, organisations
face a persistent gap between the technical execution of pipelines and the
evidentiary requirements of regulators~\cite{cps234,cps230}.

This thesis addresses that gap. Specifically, it investigates how a
\emph{DevSecOps-enabled toolchain} can enhance auditability and data governance
across the ETL lifecycle. Whereas DevSecOps has institutionalised security and
compliance for software artifacts, an equivalent model for data artifacts has yet
to emerge. The absence of such a model exposes regulated organisations to
compliance risk, weakens governance, and impairs the ability to demonstrate
trustworthy data management.

\section{Research Objectives}
This research adopts a Design Science methodology
\cite{hevner_design_2004} to design, implement, and evaluate a
DevSecOps-based toolchain that enhances auditability and data governance within
ETL workflows. The study follows an iterative process of artifact design,
prototype construction, and evaluation against explicit metrics for lineage
completeness, reproducibility, and tamper detection.

The investigation is guided by the following research questions:

\textbf{Main Research Question:}
\begin{quote}
How can a DevSecOps toolchain be designed to support ETL workflows in regulated
industries, enhancing auditability and data governance?
\end{quote}

\textbf{Sub-questions:}
\begin{enumerate}
  \item How do data workflows such as ETL manifest in regulated industries, and
  what auditability and data governance challenges do they face?
  \item How do DevSecOps practices enhance auditability and data governance
  within the software development lifecycle?
  \item In what ways do the stages of the data lifecycle align with or diverge
  from the stages of the software development lifecycle, and how does this
  affect the integration of DevSecOps practices?
  \item What are the key categories of tooling in DevSecOps and in ETL
  pipelines, and how do their features compare in terms of supporting
  auditability and data governance?
\end{enumerate}

\section{Contributions}
This thesis makes three primary contributions.  
First, it proposes an \textbf{architectural model} for a DevSecOps-enabled ETL
toolchain, aligning the data lifecycle with secure software delivery practices
to improve auditability and data governance.  
Second, it presents a \textbf{prototype implementation} that integrates
infrastructure as code, automated security scanning, and immutable logging into
an operational ETL environment based on Airflow, Spark, and AWS components.  
Third, it provides an \textbf{evaluation framework} that quantifies
auditability and data governance using traceability, reproducibility, and
tamper-resistance metrics. Together, these contributions establish a
proof-of-concept for embedding DevSecOps principles within data-centric systems
in regulated environments.

\section{Thesis Structure}
The remainder of this thesis is organised as follows:
\begin{itemize}
  \item \textbf{Chapter~2} reviews the literature on ETL pipelines, data
  governance frameworks, and DevSecOps practices.
  \item \textbf{Chapter~3} details the Design Science methodology and research
  design.
  \item \textbf{Chapter~4} describes the baseline ETL pipeline and identifies
  auditability and data governance gaps.
  \item \textbf{Chapter~5} presents the proposed DevSecOps-integrated
  architecture.
  \item \textbf{Chapter~6} outlines the prototype implementation and tooling
  configuration.
  \item \textbf{Chapter~7} evaluates the system against auditability and data
  governance metrics.
  \item \textbf{Chapter~8} discusses the implications, limitations, and avenues
  for future research.
  \item \textbf{Chapter~9} concludes with reflections on how DevSecOps can
  enable trustworthy data governance in regulated environments.
\end{itemize}